%% ====================================================================
%%  Tecnología ganadora
%% ====================================================================
\section{Tecnología ganadora}\label{sec:ganadora}

\begin{enumerate}[label=\textbf{\arabic*.}, leftmargin=*, itemsep=0.9em]
  \item \textbf{Proxy-Caché / Gateway}

  Del análisis comparativo realizado mediante la metodología DAR, se establece una arquitectura híbrida para la capa de proxy-caché del grupo de investigación GRID, combinando las fortalezas de dos tecnologías líderes:

  HAProxy como Proxy Inverso y Gestión de Certificados. HAProxy se posiciona como la tecnología ganadora para la función de proxy inverso y terminación TLS dentro del GRID, al obtener la puntuación DAR más alta del ranking (2.87/3.0). Es muy eficiente en CPU y RAM ($<$1 ms de latencia, 42k RPS, consumo CPU $\sim$50 \%), respaldada por benchmarks oficiales y una LTS garantizada hasta 2031. Para el GRID, HAProxy cumple de manera nativa y excelente las NPO críticas de automatización de certificados SSL (Let's Encrypt integrado), gobernanza del tráfico mediante ACLs granulares, rate limiting avanzado y adaptación a cambios mediante runtime API dinámica sin downtime. Su rol como proxy inverso garantiza la publicación segura de servicios internos hacia la comunidad académica, eliminando advertencias de navegadores y asegurando la confianza en los sitios web del grupo.

  NGINX como Forward Proxy y Control de Tráfico. NGINX se selecciona como la alternativa documentada y complementaria para las funciones de forward proxy, control de tráfico y aceleración en la descarga de archivos grandes. Aunque su puntuación DAR global (2.49) es inferior a HAProxy, NGINX aporta una documentación académica extensa y una caché HTTP robusta validada por F5 en 2026, aspectos relevantes para el contexto universitario del GRID. En su rol de forward proxy, NGINX permite la aceleración de descargas de ISOs y contenido académico mediante caché HTTP con cache hit ratio superior al 80 \%, mientras que sus capacidades de QoS y ACLs facilitan el control del uso no misional de la red. Su integración con la infraestructura existente del GRID es plug-and-play, minimizando la curva de aprendizaje y sirviendo como respaldo documentado ante cualquier escenario de fallo o migración futura.

  \item \textbf{Almacenamiento de Archivos de Gran Tamaño}

  Para el backend de almacenamiento de objetos grandes (ISOs, VMs, datasets) del GRID, el análisis DAR identifica una estrategia escalonada que equilibra la producción con la agilidad del prototipo funcional requerido en el trabajo de grado:

  Ceph como Backend Principal de Producción: Ceph lidera el ranking de almacenamiento con una puntuación DAR de 2.89/3.0, siendo la tecnología ganadora indiscutible para el repositorio local robusto del GRID. Su arquitectura distribuida con replicación 3x + erasure coding (EC) y auto-healing mediante el algoritmo CRUSH, garantiza la persistencia y tolerancia a fallos de los archivos grandes. Con más de 500K IOPS en NVMe (BlueStore) y respaldo enterprise de IBM/Red Hat (versiones Reef, Squid y Tentacle activas en 2026), Ceph satisface la necesidad de un repositorio local de archivos grandes, escalabilidad sin reemplazo de hardware y alta disponibilidad. Su integración con Prometheus y Grafana permite el monitoreo en tiempo real del consumo de almacenamiento y salud del cluster, alineándose con los criterios de observabilidad del proyecto.

  SeaweedFS como Backend del Prototipo Funcional: SeaweedFS se selecciona como la tecnología ganadora para el prototipo funcional del trabajo, al ofrecer un equilibrio óptimo entre funcionalidad S3-compatible y menor complejidad operativa respecto a Ceph. Con una puntuación DAR de 2.49, SeaweedFS soporta erasure coding nativo, API S3 completa y tiempos de respuesta de 2.1 ms para archivos pequeños, siendo suficiente para validar el marco operacional de proxy-caché del GRID sin requerir la curva de aprendizaje y recursos de un cluster Ceph completo. Su soporte activo por parte de Chris Lu (v3.85+) y su arquitectura liviana permiten desplegar el prototipo en la infraestructura no productiva del grupo de investigación de manera ágil, cumpliendo con el alcance del anteproyecto y generando métricas reales de throughput y cache hit ratio.

  Ceph destaca en producción a gran escala, pero su arquitectura implica una complejidad operativa excesiva para el alcance de un prototipo funcional. SeaweedFS ofrece una API S3 completa, acceso a disco O(1) y una arquitectura ligera que se despliega en minutos sobre la infraestructura del GRID, con una huella de CPU y memoria muy inferior. Además, incorpora métricas Prometheus nativas y dashboards Grafana, junto con erasure coding y replicación configurables, aportando observabilidad y tolerancia a fallos suficientes para validar el marco operacional proxy-caché sin la curva de aprendizaje ni los requisitos de hardware de un clúster Ceph. Por ello, tras reponderar la matriz DAR con criterios centrados en el prototipo, SeaweedFS supera a Ceph tanto en la valoración individual (2,83 frente a 2,70) como en la media consolidada del equipo (2,66 frente a 2,63), por lo que se selecciona como backend del prototipo.

\end{enumerate}
